\section{Nets, generators and paths}\label{sec:objects}

\motiv{\vet{A Lieb-Robinson bound is a statement about three things: where observables live,
how the dynamics is assembled from local pieces, and which chains of pieces can carry a signal
from one region to another. In this chapter we fix the three objects. The net of
local algebras records where observables live (Definition~\ref{def:net}); the local generator
family records the dynamics as a finite sum of bounded maps, each tied to a patch of sites and
each obeying three hypotheses (Definition~\ref{def:local-generators}); the paths record the
chains (Definition~\ref{def:path}).

The path-sum bound of Chapter~\ref{sec:master} \revs{bounds the commutator of an observable $a$
on a region $X$, evolved in time, with a second observable $b$, the \emph{probe}, on a region
$Y$, by a weighted count of paths from $Y$ to $X$}. A reader who knows the literature's
interactions $\Phi(Z)$ will recognise the Hamiltonian case in Example~\ref{ex:hamiltonian}.}}

\subsection{The net of local algebras}\label{subsec:net}

\vet{Throughout, $\Lambda$ is a set, the \emph{sites}, and $\cA$ a unital C*-algebra, the
\emph{ambient algebra}. The regions are all subsets of $\Lambda$.}

\begin{definition}[net of local algebras]\label{def:net}
A \emph{net} assigns to every region $Z\subseteq\Lambda$ \vetnew{a unital C*-subalgebra
$\cA(Z)\subseteq\cA$ (a norm-closed $*$-subalgebra containing $1$)}, the \emph{local algebra}
of $Z$, monotonically: $\cA(Z)\subseteq\cA(Z')$
whenever $Z\subseteq Z'$. The net is \emph{causal} if observables of disjoint regions commute:
$ab=ba$ for $a\in\cA(Z)$, $b\in\cA(Z')$ and $Z\cap Z'=\emptyset$.
\end{definition}

\vet{The net is the one of algebraic quantum field theory, an assignment of local algebras to
regions over an abstract causal structure that says which regions are spacelike separated and
which depend on which. On a lattice the structure is the simplest one: two regions are spacelike
separated exactly when they are disjoint, and a region depends on another exactly when it is
contained in it. Causality of the net for this structure is the commutation above, and it is
the only property of the structure used in these notes. \vetnew{We
require the assignment on all subsets, including infinite ones, because the proof of the
path-sum bound uses the local algebra of the complement $Y^c$ of the probe's region, which is
infinite when $\Lambda$ is infinite and $Y$ finite.} On a finite lattice this is no restriction; on the infinite lattice of
Chapter~\ref{sec:limit} the quasi-local algebra supplies $\cA(Z)$ for every $Z$.}

\subsection{The commutator seminorm}\label{subsec:commnorm}

\motiv{\vet{A bounded map of $\cA$ can only be called local if there is a way to say how much
an observable \emph{sees} a region. The right quantity is not the norm of the observable but
how far it is from commuting with the region's algebra; an observable supported far away sees
nothing, however large it is. The commutator seminorm measures this, and we phrase hypothesis
(G1) below with it.}}

\begin{definition}[commutator seminorm]\label{def:commnorm}
For a region $S$ and $c\in\cA$,
\begin{equation}\label{eq:commnorm}
  N_S(c):=\sup\{\norm{\comm cb}: b\in\cA(S),\ \norm b\le1\},
\end{equation}
where $\comm cb=cb-bc$.
\end{definition}

\begin{proposition}[properties of the seminorm]\label{prop:commnorm}
For regions $S\subseteq S'$, $c,c'\in\cA$ and \vetnew{$z\in\C$}:
\begin{enumerate}[label=(\alph*),nosep]
\item $0\le N_S(c)\le2\norm c$, and $\norm{\comm cb}\le N_S(c)$ for every $b\in\cA(S)$ with
  $\norm b\le1$;
\item $N_S(c)=0$ when $c$ commutes with every element of $\cA(S)$;
\item $N_S(c)\le N_{S'}(c)$;
\item $N_S(c+c')\le N_S(c)+N_S(c')$, \vetnew{$N_S(zc)=|z|\,N_S(c)$}, and $|N_S(c)-N_S(c')|\le2\norm{c-c'}$.
\end{enumerate}
\end{proposition}

\begin{proof}
The set in \eqref{eq:commnorm} contains $0$ (take $b=0$) and is bounded by $2\norm c$, since
$\norm{\comm cb}\le2\norm c\norm b$; this gives (a). In (b) every element of the set is $0$.
In (c) the set for $S$ is contained in the set for $S'$ because $\cA(S)\subseteq\cA(S')$. The
first two parts of (d) follow from $\comm{c+c'}b=\comm cb+\comm{c'}b$ and
\vetnew{$\comm{zc}b=z\comm cb$} by taking suprema; the last part from the first two applied to
$c=(c-c')+c'$ and $c'=(c'-c)+c$ together with (a).
\end{proof}

\subsection{Local generators}\label{subsec:generators}

\motiv{\vet{The dynamics of a lattice system is a one-parameter group or semigroup $A\mapsto
e^{tL}A$ on the observables, and $L$ is a sum of pieces, one per interaction term. For a
Hamiltonian $H=\sum_i\Phi_i$ the pieces are the derivations $\ii\comm{\Phi_i}{\cdot}$; for an
open system they are local Lindblad generators. What the bound needs from a piece is not that it
is a derivation but three properties: it acts on an observable only through the observable's
failure to commute with its own patch; it spreads supports by at most its patch; and no
sub-collection of pieces can amplify norms. Stating the hypotheses at this level is what makes
Hamiltonian and open dynamics instances of one theorem.}}

\begin{definition}[family of local generators]\label{def:local-generators}
A \emph{family of local generators} on a net is a finite index set $I$ and, for each $i\in I$,
a finite region $S_i\subseteq\Lambda$ (the \emph{support}), a bounded linear map
$L_i:\cA\to\cA$ (the \emph{generator}) and a number $\kappa_i\ge0$ (the \emph{strength}), such
that
\begin{enumerate}[label=(G\arabic*),nosep]
\item \emph{locality in the argument:} $\norm{L_ic}\le\kappa_i\,N_{S_i}(c)$ for every $c\in\cA$;
\item \emph{support closure:} $L_i\,\cA(Z)\subseteq\cA(Z\cup S_i)$ for every region $Z$;
\item \emph{contraction:} $\norm{e^{tL_J}}\le1$ for every $J\subseteq I$ and every $t\ge0$,
  where $L_J:=\sum_{i\in J}L_i$.
\end{enumerate}
The \emph{full generator} is $L:=L_I$, the \emph{total strength} $K:=\sum_{i\in I}\kappa_i$, and
the \emph{dynamics} is $A(t):=e^{tL}a$ for $a\in\cA$ and $t\ge0$. A set $S$ \emph{crosses} a
region $Y$ if $S\cap Y\ne\emptyset$ and $S\not\subseteq Y$, that is, $S$ meets $Y$ and its
complement. For a region $Y$, $I^{\partial Y}:=\{i\in I:S_i\text{ crosses }Y\}$ is the set of
terms \emph{crossing} $Y$.
\end{definition}

\vet{\vetnew{The superscript in $I^{\partial Y}$ is for now a label for the terms crossing $Y$;
the boundary $\bd Y$ as a set of sites is defined in Chapter~\ref{sec:closed}.} The exponential
is that of the Banach algebra of bounded operators on $\cA$,
$e^{tL}=\sum_{n\ge0}t^nL^n/n!$. We ask for (G3) for every sub-family because the proof
switches off the terms crossing the probe's region and evolves with the rest; the instances verify it
uniformly (unitarity for a Hamiltonian family, Chapter~\ref{sec:qudit}; for a Lindblad family
it is carried as a hypothesis, implied by complete positivity with unit preservation,
Chapter~\ref{sec:approx}), so nothing is lost by
asking for all $J$ at once, though the proof uses only the full family
and the complements $I\setminus I^{\partial Y'}$ of the regions $Y'$ that the recursion
visits.}

\begin{proposition}[sub-families]\label{prop:subfamilies}
For a family of local generators, all $J,J'\subseteq I$, every region $Y$ and all $a,c\in\cA$:
\begin{enumerate}[label=(\alph*),nosep]
\item $L_{J\cup J'}=L_J+L_{J'}$ when $J\cap J'=\emptyset$; in particular
  $L=L_{I^{\partial Y}}+L_{I\setminus I^{\partial Y}}$;
\item $\norm{e^{tL}a}\le\norm a$ for $t\ge0$;
\item $\norm{L_Jc}\le\sum_{i\in J}\kappa_i\,N_{S_i}(c)$.
\end{enumerate}
\end{proposition}

\begin{proof}
\begin{enumerate}[label=(\alph*),nosep]
\item The splitting of a finite sum over a disjoint union.
\item (G3) for $J=I$.
\item The triangle inequality and (G1) term by term.\qedhere
\end{enumerate}
\end{proof}

\begin{example}[Hamiltonian terms]\label{ex:hamiltonian}
\vet{\vetnew{In a nonzero ambient algebra, $\cA\ne0$ (so that a unitary has norm one), let}
$\Phi_i\in\cA(S_i)$ be self-adjoint, $L_i:=\ii\comm{\Phi_i}{\cdot}$ and
$\kappa_i:=\norm{\Phi_i}$. Then (G1) holds because $\norm{L_ic}=\norm{\comm{c}{\Phi_i}}
=\norm{\Phi_i}\,\norm{\comm{c}{\Phi_i/\norm{\Phi_i}}}\le\kappa_iN_{S_i}(c)$ when
$\Phi_i\ne0$, and trivially otherwise. (G2) holds because $\Phi_ic$ and $c\Phi_i$ lie in the
algebra generated by $\cA(Z)$ and $\cA(S_i)$, which is contained in $\cA(Z\cup S_i)$. (G3)
holds because $e^{tL_J}$ is conjugation by the unitary $e^{\ii tH_J}$ with
$H_J=\sum_{i\in J}\Phi_i$, an isometry.

This is the literature's setting: an interaction $\Phi$ with terms $\Phi(Z)$ on finite sets $Z$
and the Heisenberg dynamics $\tau_t(A)=e^{\ii tH}Ae^{-\ii tH}$.}
\end{example}

\subsection{Paths}\label{subsec:paths}

\motiv{\vet{A signal from the probe's region $Y$ to the observable's region $X$ must pass from
term to term, the first crossing $Y$ and each crossing the last. The combinatorial object is the sequence of terms it
passes through, and the bound will count such sequences, each weighted by the product of the
strengths of its terms. We fix two conventions here that matter later. Each term crosses the
patch of the previous one, meeting it and its complement: a term inside the current patch cannot
carry the signal out of it, as Chapter~\ref{sec:master} shows, so no term follows itself and no
term is nested in its predecessor (a term may recur later, after another one). And a path ends
at the first term that touches $X$, because the recursion stops there.
Figure~\ref{fig:hypergraph} shows the crossing condition; Example~\ref{ex:twosites} and the
chain of Chapter~\ref{sec:master} show the first-contact condition.}}

\begin{definition}[first-contact paths]\label{def:path}
Let $X,Y$ be regions and $k\ge1$. A \emph{path of length $k$ from $Y$ to $X$} is a sequence
$(i_1,\dots,i_k)\in I^k$ such that
\[
  S_{i_1}\text{ crosses }Y,\qquad S_{i_{j+1}}\text{ crosses }S_{i_j}\ \ (1\le j<k),\qquad
  S_{i_j}\cap X=\emptyset\ \ (1\le j<k),\qquad S_{i_k}\cap X\ne\emptyset .
\]
Its \emph{weight} is $w(i_1,\dots,i_k):=\prod_{j=1}^k\kappa_{i_j}$. There is no path of length
$0$. The \emph{path count} $P_k(Y,X)$ is the sum of the weights of the paths of length $k$ from
$Y$ to $X$, and $P_0(Y,X):=0$.
\end{definition}

\begin{figure}[htb]
\centering
% Figure 1: the interaction hypergraph of three terms on five sites, with Y = {1}, X = {4}.
% Three panels: a path of length two; a path of length three in which each set crosses the
% previous one; a sequence that is a union-chain but not a path. Transcription (for cheap
% agents): sites 1..5 on a line; hyperedges S1 = {1,2} (box around sites 1 and 2), S2 = {1,3}
% (an arch: two legs around sites 1 and 3 joined above site 2, so that site 2 is NOT inside S2
% and S1 is not nested in S2), S3 = {3,4} (box around sites 3 and 4); Y = {1} (probe), X = {4}
% (observable). Panel 1 highlights S2 then S3 (steps 1, 2 from above). Panel 2 highlights S1,
% then S2, then S3 (steps 1, 2, 3; each set crosses the previous one; the arrow to S1 enters
% from the left). Panel 3 highlights S2 then S1 (blue, valid steps) and S3 (red, the failing
% step), and marks the empty intersection of S1 and S3 in red on a line of its own, below the
% labels Y, X and above the panel title. Blue marks the sets of the sequence in order; red marks
% the step that fails.
\begin{tikzpicture}[x=0.9cm,y=0.9cm,
  site/.style={circle,draw,fill=white,inner sep=1.6pt,font=\scriptsize},
  edge/.style={rounded corners=5pt,draw,thick,fill=none},
  hl/.style={edge,draw=Blue,line width=1.4pt},
  lab/.style={font=\scriptsize},
  stp/.style={font=\scriptsize,Blue}]
% the arch S2 = {1,3}: legs around sites 1 and 3, joined above site 2
\newcommand{\archtwo}[1]{\draw[#1] (0.75,-0.5) -- (0.75,0.7) -- (3.25,0.7) -- (3.25,-0.5)
    -- (2.75,-0.5) -- (2.75,0.5) -- (1.25,0.5) -- (1.25,-0.5) -- cycle;}
\foreach \dx/\title in {0/{path $(S_2,S_3)$}, 6/{path $(S_1,S_2,S_3)$},
    12/{not a path: $(S_2,S_1,S_3)$}} {
  \begin{scope}[shift={(\dx,0)}]
    \draw[edge] (0.55,-0.35) rectangle (2.4,0.35);               % S1 = {1,2}
    \archtwo{edge}                                              % S2 = {1,3}
    \draw[edge] (2.55,-0.35) rectangle (4.4,0.35);               % S3 = {3,4}
    \foreach \s in {1,...,5} { \node[site] (s\s) at (\s,0) {\s}; }
    \node[lab] at (1,-0.95) {$Y$};
    \node[lab] at (4,-0.95) {$X$};
    \node[lab] at (0.3,-0.5) {$S_1$};
    \node[lab] at (0.5,0.9) {$S_2$};
    \node[lab] at (4.65,-0.5) {$S_3$};
    \node[lab,font=\footnotesize] at (2.5,-1.85) {\title};
  \end{scope}
}
% panel A: S2 then S3
\begin{scope}[shift={(0,0)}]
  \archtwo{hl}
  \draw[hl] (2.55,-0.35) rectangle (4.4,0.35);
  \draw[->,Blue,thick] (2.0,1.3) -- (2.0,0.8); \node[stp] at (2.0,1.5) {$1$};
  \draw[->,Blue,thick] (3.8,1.3) -- (3.8,0.45); \node[stp] at (3.8,1.5) {$2$};
\end{scope}
% panel B: S1, S2, S3 (each crosses the previous one)
\begin{scope}[shift={(6,0)}]
  \draw[hl] (0.55,-0.35) rectangle (2.4,0.35);
  \archtwo{hl}
  \draw[hl] (2.55,-0.35) rectangle (4.4,0.35);
  \draw[->,Blue,thick] (-0.2,0) -- (0.45,0); \node[stp] at (0.1,0.25) {$1$};
  \draw[->,Blue,thick] (2.3,1.3) -- (2.3,0.8); \node[stp] at (2.3,1.5) {$2$};
  \draw[->,Blue,thick] (3.8,1.3) -- (3.8,0.45); \node[stp] at (3.8,1.5) {$3$};
\end{scope}
% panel C: S2, S1, S3 with the gap
\begin{scope}[shift={(12,0)}]
  \archtwo{hl}
  \draw[hl] (0.55,-0.35) rectangle (2.4,0.35);
  \draw[hl,draw=Red] (2.55,-0.35) rectangle (4.4,0.35);
  \draw[->,Blue,thick] (2.0,1.3) -- (2.0,0.8); \node[stp] at (2.0,1.5) {$1$};
  \draw[->,Blue,thick] (-0.2,0) -- (0.45,0); \node[stp] at (0.1,0.25) {$2$};
  \draw[->,Red,thick] (3.8,1.3) -- (3.8,0.45); \node[stp,Red] at (3.8,1.5) {$3$};
  \node[lab,Red] at (2.5,-1.4) {$S_1\cap S_3=\emptyset$};
\end{scope}
\end{tikzpicture}

\caption{\vet{The interaction hypergraph of three terms $S_1=\{1,2\}$, $S_2=\{1,3\}$,
$S_3=\{3,4\}$ on five sites, with $Y=\{1\}$ and $X=\{4\}$. Left: the path $(S_2,S_3)$ of length
two. Middle: the path $(S_1,S_2,S_3)$ of length three; the sequence $(S_2,S_2,S_3)$ is not a
path, because $S_2$ does not cross itself. Right: the sequence $(S_2,S_1,S_3)$ is not a path,
because $S_1$ and $S_3$ do not meet, although each set meets the union of $Y$ and the earlier
ones. Blue marks the sets of the sequence in order, red the step that fails. For unit strengths
the counts are $P_1=0$ and $P_k=1$ for every $k\ge2$, the single path being the alternating
sequence $(S_2,S_1,S_2,\dots,S_3)$ or $(S_1,S_2,S_1,\dots,S_3)$ ending with $S_2,S_3$. With
consecutive overlap in place of crossing, so that a term may follow itself, the first-contact
counts would be $0,1,2,4$; dropping first contact as well gives $0,1,3,8$.}}
\label{fig:hypergraph}
\end{figure}

\begin{proposition}[the path count by recursion]\label{prop:pathcount}
For all regions $X,Y$:
\begin{enumerate}[label=(\alph*),nosep]
\item $P_1(Y,X)=\sum_{i\in I^{\partial Y},\ S_i\cap X\ne\emptyset}\kappa_i$ and, for $k\ge1$,
  $P_{k+1}(Y,X)=\sum_{i\in I^{\partial Y},\ S_i\cap X=\emptyset}\kappa_i\,P_k(S_i,X)$;
\item $0\le P_k(Y,X)\le K^k$ for every $k$ (with $K^0=1$).
\end{enumerate}
\end{proposition}

\begin{proof}
\begin{enumerate}[label=(\alph*),nosep]
\item The paths of length $1$ are the single indices $i$ whose support crosses $Y$ and meets $X$,
of weight $\kappa_i$. A sequence $(i_1,\dots,i_{k+1})$ is a path of length $k+1$ from $Y$ to $X$
exactly when $S_{i_1}$ crosses $Y$ and misses $X$, and $(i_2,\dots,i_{k+1})$ is a path of length
$k$ from $S_{i_1}$ to $X$: the condition ``$S_{i_2}$ crosses $S_{i_1}$'' is the first condition of
the shorter path with $Y$ replaced by $S_{i_1}$, and the remaining conditions are unchanged.
Splitting the sum according to $i_1$ and using $w(i_1,\dots,i_{k+1})=\kappa_{i_1}w(i_2,\dots,i_{k+1})$
gives the recursion.
\item By induction on $k$ using (a): $P_1\le\sum_i\kappa_i=K$ and
$P_{k+1}(Y,X)\le\sum_i\kappa_iK^k=K^{k+1}$.\qedhere
\end{enumerate}
\end{proof}

\begin{example}[two qubits]\label{ex:twosites}
\vet{Two sites $\Lambda=\{0,1\}$, each carrying a qubit, so \vetnew{$\cA=\Mat_2(\C)\otimes\Mat_2(\C)$} with
\vetnew{$\cA(\{0\})=\Mat_2(\C)\otimes1$, $\cA(\{1\})=1\otimes\Mat_2(\C)$}, and the Pauli matrices
$\sigma^x_j,\sigma^y_j,\sigma^z_j$ acting on site $j$. The Hamiltonian is the single term
\[
  H=\sigma^z_0\sigma^z_1,
\]
so the family has one generator, $S_1=\{0,1\}$, $L_1=\ii\comm{H}{\cdot}$ and
$\kappa_1=\norm H=1$ (Example~\ref{ex:hamiltonian}). Take the observable $a=\sigma^x_0$ on
$X=\{0\}$ and the probe $b=\sigma^x_1$ on $Y=\{1\}$. The support $S_1$ crosses $Y$ and meets
$X$, so the only path is $(1)$, of length one: $P_1=1$ and $P_k=0$ for $k\ge2$, since a path of
length two or more would have to miss $X$ at its first step.

The dynamics is explicit. Since $H$ anticommutes with $\sigma^x_0$,
$e^{\ii tH}\sigma^x_0=\sigma^x_0e^{-\ii tH}$, so
$A(t)=\sigma^x_0e^{-2\ii tH}=\sigma^x_0(\cos2t-\ii\sin2t\,\sigma^z_0\sigma^z_1)$ and
\[
  \comm{A(t)}{\sigma^x_1}=-\ii\sin2t\,\sigma^x_0\sigma^z_0\comm{\sigma^z_1}{\sigma^x_1}
  =2\sin2t\,\sigma^x_0\sigma^z_0\sigma^y_1,
\]
of norm $2|\sin2t|$. The path-sum bound of Chapter~\ref{sec:master} will give $4t$.}
\end{example}

\subsection*{Notes and references}
\vet{\vetnew{The net of local algebras is the one of algebraic quantum field theory, restricted to a
lattice. The bound goes back to \cite{lieb-robinson}; its form for interactions on general
lattices, decaying with the distance, is that of \cite{hastings-koma,nachtergaele-sims}.} The
commutator seminorm is the quantity $\sup_{\norm B=1}\norm{\comm AB}$ that the review
\cite{chen-lucas-yin} uses for local approximants (Proposition~4.1 there), here promoted to the
hypothesis (G1). The literature states the dynamics through an interaction $\Phi$
(\cite{nachtergaele-sims}, eqs.~(2.5)--(2.6); \cite{nsy2019}, \S2) or through edge Hamiltonians
$H_e$ (\cite{chen-lucas-yin}, \S3); both are instances of Definition~\ref{def:local-generators}
(Chapters~\ref{sec:qudit} and~\ref{sec:approx}). The crossing condition is that of the chains
of \cite{nsy2019}, eqs.~(3.41)--(3.42), which run over sets meeting the previous set and its
complement (their chains start at the observable's region and ours at the probe's;
Chapter~\ref{sec:master} reconciles the orientation); the sequences of \cite{chen-lucas-yin},
Theorem~3.4, eq.~(3.34), run from a set meeting $Y$ to a set meeting $X$ with consecutive overlap and distinct
consecutive terms, a larger class; the first-contact condition is ours. \emph{Union-chains}, sequences
in which each set meets the union of $Y$ and the earlier sets, are a tempting but wrong
bookkeeping object: the three-term example of Figure~\ref{fig:hypergraph} separates the two
notions, and there are four-term hypergraphs on three sites for which union-chains are not
dominated by the site-matrix bound of Chapter~\ref{sec:closed}, whereas paths always are.}

\paragraph{Checks.} {\small The counts of Figure~\ref{fig:hypergraph} and of
Example~\ref{ex:twosites} are computed by \texttt{checks/m1a\_check.py} and
\texttt{checks/path\_check.py}; the four-term hypergraph is in the latter's section~2.}
